% Generated draft: requires mathematical and source review.
\documentclass[11pt]{article}
\usepackage[margin=25mm]{geometry}
\usepackage{fontspec}
\setmainfont{Latin Modern Roman}
\usepackage{amsmath,amssymb,mathtools,unicode-math}
\setmathfont{Latin Modern Math}
\usepackage{enumitem,needspace,xurl,hyperref}
\setlength{\parindent}{0pt}
\setlength{\parskip}{5pt}
\setlength{\emergencystretch}{3em}
\newcommand{\N}{\mathbb{N}}
\newcommand{\Z}{\mathbb{Z}}
\newcommand{\Q}{\mathbb{Q}}
\newcommand{\R}{\mathbb{R}}
\newcommand{\C}{\mathbb{C}}
\newcommand{\F}{\mathbb{F}}
\newcommand{\A}{\mathbb{A}}
\newcommand{\PP}{\mathbb{P}}
\newcommand{\E}{\mathbb{E}}
\newcommand{\eps}{\varepsilon}
\newcommand{\dd}{\,\mathrm{d}}
\begin{document}
\section*{UnsolvedMath ID: 20000229}
\textbf{Source problem code:} catalogue entry
\subsection{Definitions and mathematical statement}
Let \(X\) be a general smooth projective curve of genus \(g\), and let \(L\) be a general degree-\(d\) line bundle with \(h^0(X,L)=r+1\), \(r\geq3\). Put
\[
\rho(g,r,d)=g-(r+1)(g-d+r)\geq0,
\qquad\mu_m:\operatorname{Sym}^mH^0(X,L)\longrightarrow H^0(X,L^{\otimes m}).
\]
Maximal rank means
\[
\operatorname{rank}\mu_m=\min\!\left\{\binom{r+m}{m},h^0(X,L^{\otimes m})\right\},
\]
equivalently injectivity or surjectivity. The source asks to advance a tropical-independence proof from quadrics \(m=2\) to cubics \(m=3\), rather than merely assert this rank equality. Its formulation using all of \(H^0(X,L)\) implicitly uses a complete linear series; if an arbitrary Brill--Noether \(g_d^r\) is intended, a chosen \((r+1)\)-dimensional subspace \(V\subset H^0(X,L)\) must replace it. The condition \(\rho\geq0\) alone does not force \(h^0(L)=r+1\).

For tropicalized sections \(\varphi_i\), tropical independence of a finite collection means that no constants can be added so their pointwise minimum is attained at least twice everywhere. Cubic products correspond to sums \(\varphi_i+\varphi_j+\varphi_k\), allowing repeated indices. A precise tropical target needs the metric graph, divisor, sections, parameter range and sufficient independence count, none of which are supplied in the original record.

The attached aid reports that Larson's theorem solves the classical cubic maximal-rank assertion for general Brill--Noether curves over \(\C\), while an all-parameter tropical proof remains only partially established. It gives the following characteristic-free reduction. For \(0\ne s\in V=H^0(X,L)\), let
\[
Q_s=H^0(X,L^3)/sH^0(X,L^2),\qquad
\overline\mu_3:\operatorname{Sym}^3(V/ks)\longrightarrow Q_s.
\]
The reported snake sequence is
\[
0\to\ker\mu_2\to\ker\mu_3\to\ker\overline\mu_3
\xrightarrow{\delta}\operatorname{coker}\mu_2
\to\operatorname{coker}\mu_3\to\operatorname{coker}\overline\mu_3\to0.
\]
Thus quadratic surjectivity reduces cubic surjectivity to residual surjectivity, while quadratic injectivity reduces cubic injectivity to injectivity of \(\delta\). The aid also reports a Brill--Noether dimension orientation giving
\[
d\leq\frac{r(r+2)}2\leq\binom{r+2}{3}
\]
in its quadratic-surjective range; the full hypotheses of that range are not stated here.

For \((g,r,d)=(5,3,8)\), it reports the equality-case criterion that two residual cubic relations map isomorphically to the two-dimensional quadratic cokernel. This is a condition for maximal rank, not a conclusion from dimension equality alone. The aid further reports characteristic-free injective, surjective and obstructed monomial/tropical test families, without their explicit data. These reductions do not themselves supply the requested general tropical proof.
\subsection{Short English statement}
Establish cubic maximal rank by the proposed tropical-independence method for general Brill--Noether curves. The supplied record reports the classical assertion solved by Larson over \(\C\), while distinguishing the still-partial tropical route and giving a section-quotient criterion that relates cubic rank to quadratic rank and a connecting map.
\subsection{Sources}
\begin{itemize}
\item Local UnsolvedMath registry, permanent ID \(20000229\), exact public code \texttt{catalogue entry}; canonical location \texttt{aim-algebraic-geometry-notes.json}, notes[228].
\item AIM workshop ``Degenerations in algebraic geometry,'' September 2016, Problem 8, suggested by David Jensen; supplied URL \url{https://aimath.org/pastworkshops/degenalggeomproblems.pdf}. The PDF was not consulted in this drafting step. The original literature field is empty.
\item Attached research attempt 1, ``Cubic maximal rank after the classical theorem,'' attributing the classical theorem to Larson without a full publication citation, and reporting tropical test families without their data.
\end{itemize}
\end{document}
